\par\medskip\noindent
2.1.9. La definición dada en (2.1.1) no es simétrica. Intercambiando
los papeles de $P$ y $Q$ en la construcción precedente, puede
definirse un homomorfismo análogo
\begin{equation}\tag{2.1.9.1}
 \beta^n_{P,Q}:{}_nP\otimes{}_nQ\longrightarrow\uTor_1(P,Q),
\end{equation}
reduciéndose primero al caso $nP=0$ y utilizando después una
resolución plana \underline{de} $Q$. Equivalentemente, puede definirse
$\beta$ en términos de $\alpha$ exigiendo la conmutatividad del
diagrama siguiente:
\begin{equation}\tag{2.1.9.2}
\begin{tikzcd}[column sep=huge,row sep=huge]
 {}_nP\otimes{}_nQ
   \arrow[r,"\sim"]
   \arrow[d,"{\beta^n_{P,Q}}"']&
 {}_nQ\otimes{}_nP
   \arrow[d,"{\alpha^n_{Q,P}}"]\\
 \uTor_1(P,Q)\arrow[r,"\sim"]&
 \uTor_1(Q,P).
\end{tikzcd}
\end{equation}
Las flechas horizontales son los isomorfismos de simetría. Con esta
convención, la relación entre $\alpha$ y $\beta$ es
\begin{equation}\tag{2.1.9.3}
 \beta^n_{P,Q}=-\alpha^n_{P,Q}.
\end{equation}
En otros términos:

\medskip
\noindent\underline{Lema 2.1.10.}
\underline{Sean $P$ y $Q$ dos grupos conmutativos y sea $n>0$ un
entero. Entonces el diagrama siguiente es anticonmutativo, donde
$\alpha$ está definido en (2.1.1):}
\begin{equation}\tag{2.1.10.1}
\begin{tikzcd}[column sep=huge,row sep=huge]
 {}_nP\otimes{}_nQ
   \arrow[r,"\sim"]
   \arrow[d,"{\alpha^n_{P,Q}}"']&
 {}_nQ\otimes{}_nP
   \arrow[d,"{\alpha^n_{Q,P}}"]\\
 \uTor_1(P,Q)
   \arrow[r,"\sim"]
   \arrow[ur,phantom,"-" description]&
 \uTor_1(Q,P).
\end{tikzcd}
\end{equation}
\underline{Las flechas horizontales son los isomorfismos de simetría.}

%%%%%%%%%%  scan idx 242  |  folio 231  %%%%%%%%%%

\medskip
\noindent\underline{Demostración.}
El procedimiento de 2.1.8 nos reduce inmediatamente al caso
$P=Q=(\mathbb Z/n\mathbb Z)_T$. Hay que demostrar entonces que los
dos isomorfismos (2.1.8.1), obtenidos resolviendo mediante (2.1.8.2)
el primer y el segundo argumento de $\uTor_1$, respectivamente,
difieren en un signo (equivalentemente, que el isomorfismo de simetría
en $\uTor^{\mathbb Z}_1(\mathbb Z/n\mathbb Z,\mathbb Z/n\mathbb Z)$
es $-\id$). Podemos suponer evidentemente que $T$ es el topos puntual,
de modo que trabajamos con grupos conmutativos ordinarios. Sea
$L_\bullet$ el complejo
\[
 L_\bullet=\bigl(0\longrightarrow\mathbb Z
 \xrightarrow{n}\mathbb Z\longrightarrow0\bigr),
\]
considerado como una resolución plana de
$P=Q=\mathbb Z/n\mathbb Z$. Utilicemos
$L_\bullet\otimes L_\bullet$ para calcular
\[
 \Tor_1(P,Q)=\Tor_1(\mathbb Z/n\mathbb Z,\mathbb Z/n\mathbb Z).
\]
Los isomorfismos con
$H_1(L_\bullet\otimes\mathbb Z/n\mathbb Z)$, que define $\alpha$,
y con $H_1(\mathbb Z/n\mathbb Z\otimes L_\bullet)$, que define
$\beta$, están inducidos, respectivamente, por las aumentaciones
$L_\bullet\longrightarrow\mathbb Z/n\mathbb Z$ de los factores
segundo y primero de $L_\bullet\otimes L_\bullet$. Se sigue
inmediatamente que $\alpha(1,1)$ está representado por el ciclo
\[
 -1_0\otimes1_1+1_1\otimes1_0
 \in(L_\bullet\otimes L_\bullet)_1
 =L_0\otimes L_1+L_1\otimes L_0,
\]
donde $1_0$ y $1_1$ denotan el elemento $1$ de $L_0$ y $L_1$,
respectivamente, mientras que $\beta(1,1)$ está representado por
$1_0\otimes1_1-1_1\otimes1_0$. Esto prueba 2.1.10.

Para terminar estas generalidades sobre los acoplamientos $\alpha$,
queda por precisar las relaciones entre los homomorfismos
$\alpha^n_{P,Q}$ cuando varía $n$. El resultado principal es el
siguiente.

\medskip
\noindent\underline{Lema 2.1.11.}
\underline{Sean $P,Q$ dos grupos conmutativos de $T$, y sean $n,n'$
enteros positivos tales que $n'$ sea múltiplo de $n$, digamos
$n'=mn$. Entonces el diagrama siguiente es conmutativo:}

%%%%%%%%%%  scan idx 243  |  folio 232  %%%%%%%%%%

\begin{equation}\tag{2.1.11.1}
\begin{tikzcd}[column sep=large,row sep=large]
 &{}_{n'}P\otimes{}_nQ
   \arrow[ld,"P_m\otimes\id_{{}_nQ}"']
   \arrow[rd,"\id_{{}_{n'}P}\otimes i_m"]&\\
 {}_nP\otimes{}_nQ
   \arrow[rd,"{\alpha^n_{P,Q}}"']&&
 {}_{n'}P\otimes{}_{n'}Q
   \arrow[ld,"{\alpha^{n'}_{P,Q}}"]\\
 &\uTor_1(P,Q)&
\end{tikzcd}
\end{equation}
\underline{Aquí $P_m:{}_{n'}P\longrightarrow{}_nP$ es el
homomorfismo inducido por $m\cdot\id_P$, e
$i_m:{}_nQ\longrightarrow{}_{n'}Q$ es la inclusión. Asimismo, es
conmutativo el diagrama \sisi{(2.11.2)}{(2.1.11.2)}, simétrico de
\sisi{(2.11.1)}{(2.1.11.1)}, cuyos vértices son los mismos que los de
\sisi{(2.11.1)}{(2.1.11.1)}, salvo el vértice superior
$\,{}_{n'}P\otimes{}_nQ$, que se reemplaza por
$\,{}_nP\otimes{}_{n'}Q$.}

\medskip
Basta demostrar la primera afirmación; la segunda se deduce por
simetría en virtud de 2.1.10. Mediante el procedimiento de 2.1.8,
quedamos reducidos al topos puntual, donde $P$ y $Q$ son grupos
conmutativos ordinarios y
\[
 P=\mathbb Z/n'\mathbb Z,\qquad Q=\mathbb Z/n'\mathbb Z.
\]
De hecho, para $P$ solo utilizamos la hipótesis $n'P=0$. Utilizando la
resolución
\[
 0\longrightarrow\mathbb Z\xrightarrow{n'}\mathbb Z
 \longrightarrow0
\]
de $Q$, identifiquemos $\Tor_1(P,Q)$ con $\,{}_{n'}P=P$. Asimismo,
utilizando
\[
 0\longrightarrow\mathbb Z\xrightarrow n\mathbb Z
 \longrightarrow0
\]
para resolver $\,{}_nQ\simeq\mathbb Z/n\mathbb Z$, identifiquemos
$\Tor_1({}_nP,{}_nQ)$ con $\,{}_nP$. La aumentación de esta última
resolución envía $1_0$ en $m\cdot1_Q$, donde $1_Q$ denota el elemento
$1$ de $Q=\mathbb Z/n'\mathbb Z$. Bajo estas identificaciones, el
homomorfismo
\[
\begin{tikzcd}[column sep=large,row sep=large]
 \Tor_1({}_nP,{}_nQ)\arrow[r]\arrow[d,no head,"\wr"']&
 \Tor_1(P,Q)\arrow[d,no head,"\wr"']\\
 {}_nP&{}_{n'}P=P
\end{tikzcd}
\]

%%%%%%%%%%  scan idx 244  |  folio 233  %%%%%%%%%%

inducido por las inclusiones $\,{}_nP\hookrightarrow P$ y
$\,{}_nQ\hookrightarrow Q$ no es sino la inclusión natural
$\,{}_nP\hookrightarrow{}_{n'}P$.

Por consiguiente, identifiquemos el vértice superior de (2.1.11.1)
con $P/nP$ mediante el generador $m\cdot1_Q$ de $\,{}_nQ$, que
identifica $\,{}_nQ$ con $\mathbb Z/n\mathbb Z$. El diagrama se
convierte en
\[
\begin{tikzcd}[column sep=large,row sep=large]
 &P/nP
   \arrow[ld,"x\longmapsto mx"']
   \arrow[rd,"x\longmapsto mx"]&\\
 {}_nP\arrow[rd,"x\longmapsto x"']&&
 P\arrow[ld,"x\longmapsto x"]\\
 &{}_{n'}P&
\end{tikzcd}
\]
que es manifiestamente conmutativo.

\medskip
\noindent\underline{Corolario 2.1.12.}
\underline{Con la notación de 2.1.11, el diagrama siguiente es
conmutativo:}
\begin{equation}\tag{2.1.12.1}
\begin{tikzcd}[column sep=huge,row sep=huge]
 {}_{n'}P\otimes{}_{n'}Q
   \arrow[r,"m\otimes m"]
   \arrow[d,"{\alpha^{n'}_{P,Q}}"']&
 {}_nP\otimes{}_nQ
   \arrow[d,"{\alpha^n_{P,Q}}"]\\
 {}_{n'}\uTor_1(P,Q)
   \arrow[r,"m"]&
 {}_n\uTor_1(P,Q).
\end{tikzcd}
\end{equation}
\underline{En la descripción de las flechas horizontales, $m$ denota
los homomorfismos $\,{}_{n'}E\longrightarrow{}_nE$ inducidos por
$m\cdot\id_E$, para $E=P,Q$ y $\uTor_1(P,Q)$, respectivamente.}

\medskip
En efecto, por 2.1.11, para $x\in{}_{n'}P(S)$ e
$y\in{}_{n'}Q(S)$ se tiene
\[
 \alpha^n(mx\otimes my)=\alpha^{n'}(x\otimes my).
\]
El miembro derecho es igual a $m\alpha^{n'}(x\otimes y)$, puesto que
$x\otimes my=m(x\otimes y)$. Esto prueba la conmutatividad de
(2.1.12.1).
